%! Vector Spaces and Subspaces — Unit 3, Lesson 3.1 — Slides (Beamer)
\documentclass[aspectratio=169, 11pt]{beamer}
\usepackage{linalg-beamer}   % provides \burgundyheader{} and \sectionlabel[color]{}
\newcommand{\vv}{\mathbf{v}}
\newcommand{\ww}{\mathbf{w}}
\newcommand{\xx}{\mathbf{x}}
\newcommand{\yy}{\mathbf{y}}
\newcommand{\bb}{\mathbf{b}}
\newcommand{\zero}{\mathbf{0}}

\begin{document}

% TITLE SLIDE (CourseName is not defined in beamer, so the name is literal)
{
\setbeamercolor{background canvas}{bg=burgundy}
\begin{frame}[plain]
  \vspace{2.8cm}\hspace{0.55cm}%
  \begin{minipage}{0.9\textwidth}
    {\color{goldacc}\bfseries\small LESSON 3.1}\\[0.5cm]
    {\color{white}\bfseries\LARGE Vector Spaces and Subspaces}\\[1.0cm]
    {\color{blushmid}\small Linear Algebra~$\cdot$~Unit 3: The Four Fundamental Subspaces}
  \end{minipage}
\end{frame}
}

% HOOK
\begin{frame}[t, plain]
  \burgundyheader{Prove it: is a reachable set really closed?}
  \sectionlabel{Hook}
  \begin{columns}[T]
    \begin{column}{0.60\textwidth}
      A robot arm parks its tip at $A\xx$; the set of \emph{all} reachable tips is $C(A)$.\\[6pt]
      Say it reaches $\bb_1=A\xx_1$ and $\bb_2=A\xx_2$.\\[8pt]
      \textbf{Can it reach $\bb_1+\bb_2$?}\\[4pt]
      Use $\xx=\xx_1+\xx_2$: \[ A(\xx_1+\xx_2)=A\xx_1+A\xx_2=\bb_1+\bb_2. \]
    \end{column}
    \begin{column}{0.38\textwidth}
      \centering
      \begin{tikzpicture}[scale=0.5]
        \draw[step=1, linegray!45, thin] (-0.3,-0.3) grid (3.3,3.3);
        \draw[->, thick, charcoal] (-0.3,0)--(3.5,0);
        \draw[->, thick, charcoal] (0,-0.3)--(0,3.5);
        \draw[->, very thick, burgundy] (0,0)--(2,1) node[right]{\scriptsize $\bb_1$};
        \draw[->, very thick, royalblue] (0,0)--(1,2) node[above]{\scriptsize $\bb_2$};
        \draw[->, very thick, goldacc] (0,0)--(3,3) node[above right]{\scriptsize $\bb_1+\bb_2$};
      \end{tikzpicture}
    \end{column}
  \end{columns}
  \vspace{2pt}
  \begin{block}{Today}
    Closure is not luck --- it follows from $A(\xx+\yy)=A\xx+A\yy$.
  \end{block}
\end{frame}

% WHAT IS A VECTOR SPACE
\begin{frame}[t, plain]
  \burgundyheader{What is a vector space?}
  \sectionlabel{Name the object}
  A \textbf{vector space} is a set you can \emph{add} and \emph{scale} without ever leaving it,
  obeying the usual rules:
  \begin{itemize}
    \setlength{\itemsep}{4pt}
    \item $\vv+\ww=\ww+\vv$;\; there is a $\zero$;\; every $\vv$ has $-\vv$;\; $1\vv=\vv$;\; $c(\vv+\ww)=c\vv+c\ww$.
  \end{itemize}
  \vspace{4pt}
  \begin{block}{The model example}
    $\mathbb{R}^n$ = all vectors with $n$ real components. Every rule just says arithmetic in
    $\mathbb{R}^n$ behaves normally.
  \end{block}
\end{frame}

% SUBSPACE REQUIREMENT
\begin{frame}[t, plain]
  \burgundyheader{Subspaces, made precise}
  \sectionlabel[goldacc]{The requirement}
  A \textbf{subspace} is a subset that is \emph{itself} a vector space. The usual rules are
  inherited --- you only check the two a subset could break:
  \begin{itemize}
    \setlength{\itemsep}{5pt}
    \item[\textbf{(a)}] $\vv,\ww \in S \;\Rightarrow\; \vv+\ww \in S$ \quad(closed under \textbf{addition})
    \item[\textbf{(b)}] $\vv \in S \;\Rightarrow\; c\vv \in S$ for every $c$ \quad(closed under \textbf{scaling})
  \end{itemize}
  \vspace{4pt}
  \begin{block}{Careful}
    Rule (b) with $c=0$ forces $\zero \in S$. So ``contains $\zero$'' is \textbf{necessary but not
    sufficient} --- you must still check both rules.
  \end{block}
\end{frame}

% CATALOG + SNEAKY NON-EXAMPLE
\begin{frame}[t, plain]
  \burgundyheader{The catalog --- and a sneaky ``NO''}
  \sectionlabel{Yes vs.\ No}
  \begin{columns}[T]
    \begin{column}{0.52\textwidth}
      \textbf{Subspaces of $\mathbb{R}^2$:} $\{\zero\}$, any line through $\zero$, all of $\mathbb{R}^2$.\\[4pt]
      \textbf{Of $\mathbb{R}^3$:} also every \emph{plane through $\zero$}.\\[10pt]
      The \textbf{union of the two axes} contains $\zero$ and is closed under scaling --- but
      $\vv+\ww$ escapes, so it is \textbf{not} a subspace.
    \end{column}
    \begin{column}{0.46\textwidth}
      \centering
      \begin{tikzpicture}[scale=0.62]
        \draw[step=1, linegray!45, thin] (-2.2,-2.2) grid (2.7,2.7);
        \draw[very thick, burgundy] (-2.2,0)--(2.7,0);
        \draw[very thick, burgundy] (0,-2.2)--(0,2.7);
        \draw[->, very thick, charcoal] (0,0)--(2,0) node[below right]{\scriptsize $\vv$};
        \draw[->, very thick, charcoal] (0,0)--(0,2) node[above left]{\scriptsize $\ww$};
        \fill[royalblue] (2,2) circle (2.5pt);
        \draw[->, thick, royalblue, dashed] (2,0)--(2,2);
        \draw[->, thick, royalblue, dashed] (0,2)--(2,2);
        \node[royalblue, font=\scriptsize\bfseries, right] at (2.05,2.3){$\vv+\ww$};
        \fill[black] (0,0) circle (2.5pt);
      \end{tikzpicture}
    \end{column}
  \end{columns}
\end{frame}

% C(A) PROOF
\begin{frame}[t, plain]
  \burgundyheader{$C(A)$ is a subspace --- proof}
  \sectionlabel[goldacc]{Column space}
  Members of $C(A)$ are reachable targets $A\xx$. Take $\bb_1=A\xx_1$, $\bb_2=A\xx_2$:
  \begin{itemize}
    \setlength{\itemsep}{6pt}
    \item \textbf{Add:} $\bb_1+\bb_2=A\xx_1+A\xx_2=A(\xx_1+\xx_2)$ --- reachable.
    \item \textbf{Scale:} $c\,\bb_1=c(A\xx_1)=A(c\xx_1)$ --- reachable.
  \end{itemize}
  \vspace{4pt}
  \begin{block}{So}
    Both rules hold: $C(A)$ is a subspace of $\mathbb{R}^m$ ($m$ = rows). The precise version of
    3.0's ``combinations of reachable vectors are reachable.''
  \end{block}
\end{frame}

% ROADMAP
\begin{frame}[t, plain]
  \burgundyheader{Where this goes}
  \sectionlabel{Connections}
  \textbf{Today:} vector spaces, the subspace requirement, the catalog, and $C(A)$ proved to be a subspace.\\[8pt]
  \begin{itemize}
    \setlength{\itemsep}{4pt}
    \item \textbf{3.2} The Nullspace $N(A)$: all solutions of $A\xx=\zero$ --- the second fundamental subspace.
    \item \textbf{3.3} The Complete Solution to $A\xx=\bb$.
    \item \textbf{3.4} Independence, Basis, and Dimension.
    \item \textbf{3.5} Dimensions of the Four Subspaces.
  \end{itemize}
\end{frame}

\end{document}
